\section{Theory}
\label{sec:theory}

Let \(\mathcal A\) and \(\mathcal U\) denote Banach spaces of input and output fields. A PDE family induces a solution operator \(\mathcal S:\mathcal A\to\mathcal U\). Let \(G\) be a Lie group acting on the full problem family through transformations \(T_g^{\mathcal A}\) and \(T_g^{\mathcal U}\). Closure of the full problem class under \(G\) means that the equation, coefficients, forcing terms, boundary or initial conditions, domain class, and observation map transform consistently. Under this assumption the exact solution operator is equivariant: \(\mathcal S(T_g^{\mathcal A}a)=T_g^{\mathcal U}\mathcal S(a)\). For a learned operator \(\mathcal G_\theta\), define the operator-level equivariance defect \(\Delta_\theta(a,g)=\mathcal G_\theta(T_g^{\mathcal A}a)-T_g^{\mathcal U}\mathcal G_\theta(a)\). Our regularizer jointly samples a Lie-algebra direction \(X\), a step \(\epsilon\), and \(g=\exp(\epsilon X)\), then penalizes
\begin{equation}
\mathcal L_{\mathrm{orb}}(\theta)=\mathbb E_{a,X,\epsilon}\left[\frac{\|\Delta_\theta(a,\exp(\epsilon X))\|_{\mathcal U}^2}{\epsilon^2+\eta}\right].
\end{equation}
For gridded experiments, \(\|\cdot\|_{\mathcal U}^2\) is discretized as the
mean square over spatial and output entries, matching the implementation in
Section~\ref{sec:method}.
The objective is label-free because both terms are model predictions transformed by known group actions. For the N64 square sampling law, \(X=\delta/\|\delta\|_2\) away from the measure-zero identity and \(\epsilon=\|\delta\|_2\); direction and step are therefore jointly induced rather than independently sampled.

\paragraph{Infinitesimal interpretation.} Assume \(T_{\exp(\epsilon X)}^{\mathcal A}\), \(T_{\exp(\epsilon X)}^{\mathcal U}\), and \(\mathcal G_\theta\) are differentiable at \(\epsilon=0\). Let \(X^{\mathcal A}a\) and \(X^{\mathcal U}u\) denote the induced infinitesimal actions. Taylor expansion gives
\begin{equation}
\mathcal G_\theta(T_{\exp(\epsilon X)}^{\mathcal A}a)-T_{\exp(\epsilon X)}^{\mathcal U}\mathcal G_\theta(a)=\epsilon\left(D\mathcal G_\theta(a)[X^{\mathcal A}a]-X^{\mathcal U}\mathcal G_\theta(a)\right)+O(\epsilon^2).
\end{equation}
Thus \(\|\Delta_\theta(a,\exp(\epsilon X))\|_{\mathcal U}^2/\epsilon^2\) converges to \(\|D\mathcal G_\theta(a)[X^{\mathcal A}a]-X^{\mathcal U}\mathcal G_\theta(a)\|_{\mathcal U}^2\) as \(\epsilon\to0\). For the implemented denominator \(\epsilon^2+\eta\), this tangent interpretation applies when \(\epsilon^2\gg\eta\), or along a limit with \(\eta=o(\epsilon^2)\); with fixed positive \(\eta\), the loss instead vanishes as \(\epsilon\to0\). LOCO is therefore a finite-difference-based regularizer that avoids explicit Jacobian--vector products, not an unbiased tangent estimator at the finite radii used experimentally. For nonzero \(\epsilon\), the normalized defect also contains higher-order variation of the learned map and group orbit. Enforcing commutation over a sampled neighborhood can therefore constrain finite displacements that an identity-only directional derivative does not; the tangent and finite-step controls in Section~\ref{sec:experiments} test this distinction.

\paragraph{Orbit error decomposition.} If the true solution operator is equivariant and \(T_g^{\mathcal U}\) is bounded, then for any \(a\in\mathcal A\) and \(g\in G\),
\begin{align}
\|\mathcal G_\theta(T_g^{\mathcal A}a)-\mathcal S(T_g^{\mathcal A}a)\|_{\mathcal U}
&\le \|\mathcal G_\theta(T_g^{\mathcal A}a)-T_g^{\mathcal U}\mathcal G_\theta(a)\|_{\mathcal U} \nonumber\\
&\quad+\|T_g^{\mathcal U}\|\,\|\mathcal G_\theta(a)-\mathcal S(a)\|_{\mathcal U}.
\end{align}
The first term is the learned equivariance defect and the second term is the in-distribution prediction error transported by the output action. This inequality motivates reducing \(\mathcal L_{\mathrm{orb}}\) when test inputs differ from training inputs by known PDE symmetries.

\paragraph{Approximate symmetries and cost.} For partial or approximate symmetries, replace the output norm by a weighted norm \(\|v\|_w^2=\|w_{a,g}\odot v\|^2\), where \(w_{a,g}\) can encode valid overlap after transformations, boundary exclusions, observation masks, or regions where transformed coefficients and forcing terms remain physically meaningful. With one sampled transform per minibatch, the LOCO term uses one evaluation on the original input, one on the transformed input, and a known output transformation. At deployment there is no LOCO module or additional operator call: inference uses the same base graph \(\mathcal G_\theta(a)\), and the measured latency differences are within run-to-run variation.
